% ============================================================
% JMGM Methods — F3-b re-scope (2026-10-06). Numbers trace to the ledger and
% the scripts in code/.
% ============================================================

\section{Methods}\label{sec:methods}

\subsection{Experimental target and dataset}\label{para:validation_hierarchy}

The prediction target is the \emph{experimental} singlet--triplet gap
$\Delta E_\text{ST}$. We took literature-reported values from the ChemDataExtractor-generated TADF
database of Huang and Cole~\cite{HuangCole2024} and matched them, by compound name, to molecules in
our descriptor set, retaining values in the physical range
\qtyrange{-0.3}{1.5}{\electronvolt}. Where a compound carried multiple independent
reports we used the median and recorded the inter-report spread; a single report sits
\qty{0.072}{\electronvolt} from the median (RMS, weighting each molecule equally; the
report-pooled value is \qty{0.0905}{\electronvolt}). For the \num{102} multi-report
molecules the median carries an RMS standard error of \qty{0.049}{\electronvolt}; with the
\num{129} single-report labels counted at single-report scatter, the label noise over all
\num{231} molecules is about \qty{0.072}{\electronvolt} RMS (\cref{subsec:ceilings}). The final benchmark comprises \num{231} donor--acceptor
molecules spanning \num{212} distinct Bemis--Murcko scaffolds~\cite{Bemis1996}.

Each molecule was assigned one \emph{primary DOI}: the source publication that contributes
the most reports for that molecule (ties are broken by the
code's report-count ordering). This gives \num{173} primary-DOI groups for the \num{231} molecules, of
which \num{128} contain one molecule and \num{45} contain more than one
(\texttt{code/jmgm\_doi\_group\_ranking\_check.py}).

The corpus does not record measurement medium (solvent, host) for individual
values, so the per-molecule median is a consensus over uncontrolled conditions
rather than a single, solvent-specific state. Because environmental electrostatics
and conformational disorder are recognised as primary obstacles to accurate TADF
modelling,\cite{Olivier2017} this consensus is best read as the central tendency of
experimentally accessible $S_1$--$T_1$ gaps; it trades physical specificity for
reduced variance. For the \num{102} molecules with repeat reports, a single random
report deviates from this median by \qty{0.072}{\electronvolt} on
average---comparable to the model error itself
(Table~\ref{SI-tab:solvent}). For those molecules the median is known more precisely,
to \qty{0.049}{\electronvolt} RMS standard error.

The sTDA $S_1$ and $T_1$ energies are excluded from the NTO set because supplying them
reproduces a computed gap exactly: under scaffold cross-validation a linear model given
both energies recovers $E_{S_1}-E_{T_1}$ with $R^2 = 1.000$, and a random forest reaches
$R^2 = 0.59$ on the same features (Section~\ref{SI-sec:leakage}).

\subsection{Cheap electronic-structure inputs and physics-informed descriptors}

Excited-state quantities used to build descriptors were obtained at the
semi-empirical level: GFN2-xTB~\cite{Bannwarth2019} geometries
(\texttt{xtb} version 6.7.1) followed by sTDA~\cite{Grimme2017} for the lowest singlet
and triplet states. From the natural transition orbitals
(NTOs)~\cite{Martin2003} of each excitation we computed spatial overlap
descriptors that encode exchange-integral physics:
\begin{align}
  S_\text{he} &= \int \phi_\text{hole}(\mathbf{r})\,\phi_\text{electron}(\mathbf{r})\,d\mathbf{r}, \label{eq:overlap}\\
  \Delta\lambda_\text{A} &= |\lambda_\text{A}(S_1) - \lambda_\text{A}(T_1)|, \label{eq:delta_lambda}\\
  \text{Char}_\text{diff}^2 &= \Big[\textstyle\sum_i |q_i(S_1) - q_i(T_1)|\Big]^2, \label{eq:char_diff}\\
  \Delta S_\text{he} &= |S_\text{he}(S_1) - S_\text{he}(T_1)|. \label{eq:delta_overlap}
\end{align}
$S_\text{he}$ quantifies hole--electron overlap (related to the exchange
integral), $\Delta\lambda_\text{A}$ acceptor-participation differences,
$\text{Char}_\text{diff}^2$ charge-transfer character changes, and
$\Delta S_\text{he}$ the singlet--triplet overlap difference. The full set of
\num{35} NTO descriptors (the $S_1$/$T_1$ energy scalars excluded) is listed in
the Supporting Information; these descriptors require the sTDA excited-state
calculation, even though they omit the energies themselves.
The ground-state wavefunction for sTDA was generated with \texttt{xtb4stda}, and
excitations were computed with \texttt{stda -xtb -e 10} (sTDA-xTB parametrisation,
\qty{10}{\electronvolt} excitation-energy window; \texttt{-t} for triplets), all in the
gas phase; descriptors use the sTDA rather than the sTD-DFT (\texttt{-rpa}) results.

\subsection{Feature sets and machine-learning models}

We compared three feature representations on the same molecules. Two require no
quantum chemistry, being computed from the 2D structure alone:
(i)~RDKit physicochemical descriptors~\cite{Landrum2016} and (ii)~Morgan
fingerprints (radius~2, \num{2048} bits; RDKit
\texttt{GetMorganFingerprintAsBitVect} with its defaults, i.e.\ no chirality and
no feature invariants). The third, (iii)~the NTO spatial
descriptors above, is derived from the sTDA excited state but excludes the energy
scalars. A reference set that additionally includes the sTDA $S_1$/$T_1$
energies was also evaluated. We used random-forest
regressors~\cite{Breiman2001} (\num{400} trees, scikit-learn~\cite{Pedregosa2011}),
which outperformed support-vector regression, ridge regression and a mean-value
baseline on this task.

\paragraph{Evaluation.}
Error metrics, \cref{tab:feature_comparison}, the baselines of Table~\ref{SI-tab:baselines}, the SHAP analysis and the
label-noise analyses use Bemis--Murcko scaffold \texttt{GroupKFold} ($k=5$); the
source-paper (DOI-grouped) results and the holdout-unit controls below use the
holdout unit stated with each number. Because \qty{92.9}{\percent} of scaffolds are singletons, this split is
numerically close to a random split; we therefore also report five-fold source-paper
(DOI-grouped) out-of-fold predictions. That grouping holds out source publications and any synthetic
series or measurement practices correlated with them, but does not identify a
laboratory or control the protocol (\cref{subsec:benchmark}). The corpus records
no solvent, host or temperature per value, so a paper DOI is only a proxy for
shared conditions, not a measured protocol.
We benchmark against a mean-value predictor, the raw sTDA gap, and a
range-separated TD-DFT reference gap (\cref{subsec:benchmark}), and report
\qty{95}{\percent} confidence intervals from \num{2000} bootstrap resamples of the
out-of-fold predictions.
For the DOI-grouped rank contrast, we also resampled the \num{173} primary DOI
groups as clusters (\num{5000} paired percentile-bootstrap draws, seed \num{54321})
while keeping the original out-of-fold predictions fixed. As a partition
sensitivity check, we refitted both random forests in \num{20} source-paper
\texttt{GroupShuffleSplit} partitions (\qty{20}{\percent} of DOI groups held out;
seed \num{17}) and computed Spearman correlation on each test partition. These
overlapping partitions are descriptive, not independent confirmations or a
confidence interval for new laboratories.

\paragraph{Holdout-unit controls.}
To separate the effect of the holdout unit from the effect of the error-versus-ranking metric,
we compared three units on the same \num{231} molecules: single molecule (random
folds), Bemis--Murcko scaffold (\num{212} groups) and primary source publication
(\num{173} groups). For each unit we drew \num{50} random assignments of groups to
\num{5} folds, fitted both random forests (\num{400} trees) on each assignment and
computed pooled out-of-fold Spearman $\rho$ and MAE for Morgan and NTO features. The
spread over the \num{50} assignments (2.5th--97.5th percentiles) describes
fold-assignment variability; it is not a confidence interval. The \num{20}-partition
\texttt{GroupShuffleSplit} count was repeated for each unit. On a fixed scaffold split we
also resampled scaffold clusters (\num{5000} draws). The same \num{50} assignments were
reused to split the molecules into \num{103} from papers with at least one other molecule
in the corpus (multi-molecule papers) and \num{128} from single-molecule papers, and
to compute (i)~$\rho$ within each subset and (ii)~a sibling-pull statistic, the Spearman
correlation between a molecule's out-of-fold prediction and the leave-self-out mean
label of the other molecules from the same paper (\num{103} molecules). Code:
\texttt{code/a017\_split\_unit\_controls.py} and
\texttt{code/a019\_series\_shortcut\_test.py}; outputs in
\texttt{data/a017\_split\_unit\_controls.json} and
\texttt{data/a019\_series\_shortcut\_test.json}. The same \num{50} molecule- and
paper-holdout assignments were repeated with the \num{217} RDKit 2D descriptors, and
within-paper ranking was measured as the Spearman correlation between paper-centred
labels and paper-centred predictions on the \num{103} multi-paper molecules. For the
\num{128} single-paper molecules, out-of-fold predictions were averaged over the \num{50}
paper-holdout assignments and the NTO-minus-Morgan $\rho$ difference was given a paired
molecule bootstrap interval (\num{5000} draws, seed \num{54321}); code
\texttt{code/a020\_controls.py}, output \texttt{data/a020\_controls.json}. These analyses
(\texttt{a017}--\texttt{a020}) were designed after the fixed-split result
(\texttt{a016}) had been seen; the decision rule of each was fixed before it was run and
is recorded in the script's docstring. Code and full per-split values for
the fixed split are in \texttt{code/jmgm\_doi\_group\_ranking\_check.py} and
\texttt{data/a016\_jmgm\_doi\_group\_ranking.json}.

\paragraph{TD-DFT reference gap.}
For a higher-level reference we computed vertical \deltaest\ for all \num{231}
molecules at wB97X-D4/def2-SVP with RIJCOSX and the Tamm--Dancoff approximation
(ORCA~6~\cite{Neese2020}; \texttt{nroots}~8, triplets included; GFN2-xTB
ground-state geometry). All \num{231}/\num{231} jobs converged
(\qty{76.2}{\hour} elapsed on \num{8} processes, so \num{609} core-hours; median
\qty{16.1}{\minute} per molecule, slowest \qty{2.1}{\hour}). These are measured wall times
from \texttt{results/tddft\_reference\_231.csv}, not a scaling extrapolation. The raw
gap, an in-fold linearly calibrated gap, and the gap as an added descriptor in
the random forest model are evaluated against the mean-value baseline on the same
scaffold and DOI cross-validation folds (\cref{subsec:benchmark};
Table~\ref{SI-tab:tddft_headtohead}).
B3LYP/6-31G(d) and CAM-B3LYP/def2-TZVP calculations on \num{7} emitters (functional
dependence of the computed gap), and CAM-B3LYP/def2-TZVP calculations on \num{14}
emitters (adiabatic gaps and the
$T_1$--$T_2$ manifold) are described in Section~\ref{SI-sec:dft_benchmark}; the
substructure-filtered PubChem library is described in Section~\ref{SI-sec:library}.

\paragraph{Interpretability and reproducibility.}
Feature importance was quantified with exact TreeSHAP~\cite{Lundberg2017} on the
fitted forest. All scripts (\texttt{code/finalize\_model.py},
\texttt{code/rebuild\_structural\_model.py}) regenerate every reported number and
\cref{fig:model_performance}.
Software: scikit-learn~1.8.0 and RDKit~2026.03.2. Every random forest used
\num{400} trees, \texttt{random\_state}~=~0 and otherwise scikit-learn defaults.
The \num{50} fold assignments per holdout unit were drawn with NumPy generator
seeds 0, 1, \ldots, 49 (one seed per assignment,
\texttt{numpy.random.default\_rng(seed)}); the paper-cluster, scaffold-cluster
and molecule bootstraps used seed \num{54321}, and the \num{20}
\texttt{GroupShuffleSplit} partitions used seed \num{17}. Data source, splits,
features, model, hyperparameters, metrics and code location are collected in
Table~\ref{SI-tab:ml_reporting}.

\paragraph{Bootstrap confidence intervals and power analysis.}
Point estimates of the rank correlation are complemented by
\num{5000}-replicate percentile bootstrap confidence intervals on the
out-of-fold predictions (\texttt{code/r2\_ci\_power\_analysis.py}).
For the headline model the Spearman $\rho = 0.36$ ($n = 231$) has a \qty{95}{\percent}~CI of
\numrange{0.24}{0.47} ($p = \num{2.6e-8}$) and the Pearson $r = 0.50$ a \qty{95}{\percent}~CI
of \numrange{0.26}{0.68}; both exclude zero, establishing statistically
significant association despite the wide $R^2$ interval.
A power analysis for the correlation test (Fisher $z$ transform, two-sided,
$\alpha = 0.05$) shows that detecting $\rho = 0.36$ at \qty{80}{\percent} power
requires only $n \approx 60$ (\qty{90}{\percent} power: $n \approx 79$); at
$n = 231$ the achieved power is $\approx 1.0$. This bears on the ranking question only:
it establishes that the $\rho$ estimate is not sample-size-limited, and it says nothing
about the width of the $R^2$ interval, which is a different statistic and is discussed on
its own terms in \cref{subsec:benchmark}. We report Spearman~$\rho$ as the primary metric
for triage utility because triage is a ranking task, not because it is the more favourable
number; the full set of metrics and intervals is collected in
Table~\ref{SI-tab:metrics_ci}.

\subsection{Data and code availability}

Labels: Huang--Cole TADF database~\cite{HuangCole2024}. Data, code and the full
analysis ledger for the three related articles: Zenodo concept DOI 10.5281/zenodo.17436069
(latest version), mirrored on GitHub (Data\_Articles\_TADF). The holdout-analysis scripts
and outputs (A-016 to A-020) are in that archive (	exttt{ML\_reproducibility/code},
	exttt{ML\_reproducibility/data}) and are also enclosed as Supplementary Data
(	exttt{supplementary\_data/}, 	exttt{JMGM\_group\_ranking\_supplement.zip}). File list: Data and code
availability statement after the Conclusions.
